%! Populations, Samples, Parameters, and Statistics — Unit 1, Lesson 1.2 — Student Packet Cover Sheet
% PILOT SHAPE (2026-09-06): modeled on AP Statistics lesson 1.4 minus its AP Practice page.
% THREE packet rows — Warm-Up, Guided Notes & Practice, Homework. The homework is the individual
% practice, started in class, and it is SCORED, so its cell is a \blank{}, never NA.
\documentclass[12pt]{article}
\usepackage{saar-article}
\usepackage{saar-boxes}
\usepackage{ltablex}
\keepXColumns
% Measured banner: the band grows to fit the title block, so a title long enough to wrap grows
% the band rather than running out of it (the stats course's \coverbanner; saar-article does not
% carry one, so it is defined here — every cover copies this block verbatim).
\newsavebox{\bannerbox}
\newlength{\bannerht}
\newlength{\coverbannerpad}
\setlength{\coverbannerpad}{0.16in}       % breathing room above and below the text
\newcommand{\coverbanner}[2]{%
  \sbox{\bannerbox}{%
    \begin{minipage}{\dimexpr\paperwidth-1.4in\relax}%
      \centering
      {\color{white}\textbf{\LARGE \CourseName}}\\[4pt]
      {\color{frostmid}\large\textbf{#1}}\\[6pt]
      {\color{white}\textbf{\large #2}}%
    \end{minipage}}%
  \setlength{\bannerht}{\dimexpr\ht\bannerbox+\dp\bannerbox+2\coverbannerpad\relax}%
  \noindent\begin{tikzpicture}[remember picture, overlay]
    \fill[cerulean] (current page.north west)
      rectangle ([yshift=-\bannerht]current page.north east);
    \node[anchor=north, inner sep=0pt]
      at ([yshift=-\coverbannerpad]current page.north) {\usebox{\bannerbox}};
  \end{tikzpicture}\par
  % The text body starts 0.6in below the page edge (geometry top=0.6in), so reserve only what
  % the band overruns that, plus a gap under the band.
  \vspace*{\dimexpr\bannerht-0.6in+0.16in\relax}%
}

\begin{document}

% Full-bleed cerulean banner, sized to the title block.
\coverbanner{Unit 1}{Lesson 1.2 \quad Populations, Samples, Parameters, and Statistics}

\vspace{0.06in}
% NAMESTRIP: this is the ONLY component that carries the name row.
\namedateperiod
\vspace{0.18in}

% GRADUAL RELEASE: the vocabulary is taught in the guided notes, so the targets name it
% plainly. The standard codes (PS.DC.1d-e) stay in the teacher-facing plan.
\begin{learningtargetbox}
\small
\begin{enumerate}[leftmargin=1.5em, itemsep=5pt]
  \item \ldots{} name the \textbf{population} and the \textbf{sample} in a study, and give
        the sample size.
  \item \ldots{} tell a \textbf{parameter} from a \textbf{statistic} by asking \emph{who the
        number describes} --- never by what the number looks like.
  \item \ldots{} name the \textbf{constraint} --- time, cost, access, destruction, or change
        --- that keeps a study from being a \textbf{census}.
  \item \ldots{} explain why two honest samples give two different statistics for one
        unchanging parameter (\textbf{sampling variability}), and say what a sample does ---
        and does \emph{not} --- let a report claim.
\end{enumerate}
\end{learningtargetbox}

\vspace{0.14in}

\begin{tocbox}
\small
\renewcommand{\arraystretch}{1.5}
\begin{tabularx}{\linewidth}{@{} c l X r @{}}
\rowcolor{frost}
\textbf{\#} & \textbf{Component} & \textbf{Description} & \textbf{Score} \\
\midrule
1 & Warm-Up & Last night's market survey again: the percent, the prediction, and the two
             groups hiding in the paragraph & \blank{1.2cm} \\
2 & Guided Notes \& Practice & Two groups and the two kinds of numbers they produce; why
             nobody runs a census; three volunteers who got three different answers ---
             then one survey we work together & \blank{1.2cm} \\
% Row 3 is the lesson's GRADED practice, started in class. Packet pages are the default; a
% Desmos activity may replace them on a given day, decided at Close & Assign. The row and its
% score cell never change.
3 & Homework & Harbor Point Community Pool: label the groups and the numbers, name the
             constraints, and explain why two surveys disagree ---
             \textbf{scored, started in class, due the first class after two study halls} & \blank{1.2cm} \\
\midrule
  & \multicolumn{2}{r}{\textbf{Total}} & \blank{1.2cm} \\
\end{tabularx}
\end{tocbox}

\vspace{0.14in}

% Keep in Mind carries the lesson's CONTENT — the rule, the test, the trap — never its process.
\begin{remindbox}
\small
A \textbf{parameter} describes the whole \textbf{p}opulation; a \textbf{statistic} describes the
\textbf{s}ample you actually collected. To tell them apart, ask \emph{who does this number
describe?} --- never \emph{what does it look like}. A percent computed from records of
\emph{everyone} is still a parameter. \textbf{Watch out:} three honest samples from one
population give three different statistics. That is \textbf{sampling variability}, not a
mistake --- the statistic moves, the parameter does not, and no sample lands exactly on it.
\end{remindbox}

\end{document}
